\scanpage{19}
La donnée d'un polygone combinatoire $P$ à $n$ côtés équivaut à la donnée d'un $D_n$-torseur $T$ \inserted{$(P={}^{T}P_n)$}, et à $T$ on associe $\,{}^{\unsure{T'|T}}E$

L'inclusion $S(P_n)\hookrightarrow E$ (\unsure{étant|c'est} une $\deleted{\unsure{P}}D_n$-application) définit une inclusion
\[
 S(P)={}^{T}S(P_n)\hookrightarrow{}^{T}E=\operatorname{Env}(P)
\]

\noindent\mbox{-{}-{}-{}-{}-}
\[
 (e,G)\mathrel{\overset{\approx}{-}}\underline{\operatorname{Tors}}(G)
\]

% slanted
\margin{DEA}
\noindent\(\left\{\begin{minipage}{0.93\linewidth}
\uline{NB}
\diagram{\begin{tikzcd}[ampersand replacement=\&,column sep=small]
\underline{\operatorname{Hom}}(\underline{\operatorname{Tors}}(G),\unsure{C|\mathcal C})
 \arrow[r,"\approx"] \&
 \begin{gathered}\text{catégor. }\deleted{\underline{\operatorname{Hom}}(G,\mathcal E)}
 \inserted{G\text{-}(\mathcal E)}\\
 \text{des objets }X\text{ de }C\\
 \text{munis d'une opération}\\
 \text{à }\deleted{\text{droite}}\inserted{\text{gauche}}\text{ de }G\\
 G\longrightarrow\operatorname{Aut}(X)^{\mathcal E}
 \end{gathered}\\
 \deleted{\unsure{T_X|P_X}\;\unsure{d'|d}} \& {}\\
 F\arrow[r,mapsto] \& F(G_d)\\
 (P\longmapsto{}^{P}X) \& X\arrow[l,mapsto]
\end{tikzcd}}

Soit $C$ catégorie, $E\in\operatorname{Ob} C$, $C_E$ \unsure{ss/cat.|s/cat.} des objets de $C$ isom. à $E$, avec iso, $G:\operatorname{Aut}_C(E)$
\[
 C_E\xrightarrow{\approx}\operatorname{Tors}(G)
\]
\diagram{\scriptsize\begin{tikzcd}[ampersand replacement=\&,row sep=large,column sep=small]
 \underline{\operatorname{Hom}}(\operatorname{Tors}(G),\mathcal E)
 \arrow[r,"\approx","{\unsure{\alpha\,\simeq|\alpha\,?}}"']
 \arrow[d,"\simeq"] \&
 \deleted{\underline{\operatorname{Hom}}(G,\mathcal E)}\inserted{G\text{-}(\mathcal E)}
 \arrow[dl,dashed,bend left=13,"\approx"']\\
 \underline{\operatorname{Hom}}(C_E,\mathcal E)
 \arrow[ur,mapsto,bend left=7,"\approx",
 "{\Phi\mapsto\Phi_{\deleted{\unsure{E|0}}}(E)}"']
 \& X\longmapsto(F\longmapsto{}^{\operatorname{Isom}(E,F)}X\unsure{)|,})\\
 {}\arrow[ur,shorten <=1em,shorten >=1.2em] \& {}
\end{tikzcd}}

Donc : la catégorie des foncteurs de $C_E$ ds $\mathcal E$ est équivalente (de façon can. à isom unique près) à celle des $G$-objets de $\mathcal E$
\end{minipage}\right.\)

% slanted
\margin{\uline{cf} $C_0$ groupoïde connexe, $\unsure{E_0|E}\in\operatorname{Ob} C_0$, $G=\operatorname{Aut}(E_0)$ \par
$\underline{\operatorname{Hom}}(C_0,\mathcal E)\longrightarrow G\text{-}(\mathcal E)$ \par
$\Phi\mathrel{\vdash\!-}\Phi(E)$ d'équivalence de catégories.}
% archive sheet number: 18
